% Einheit 3 von 3 – Wurzelgesetze und rationale Exponenten (bank/reelle-zahlen/e3.jsonl, zusammenbau v0.1)
%% TODO Zweigzeile Teil 1: was hier gelernt wird (Satz in Schülersprache aus dem Einheitstitel) – Einheit 3
\einheitenkopf[e3]{Einheit 3 von 3 · Wurzelgesetze und rationale Exponenten}
\zweigzeile{ab Kl. 9 (am Gymnasium ab Kl. 8) · keine P10-Aufgabe}
\setcounter{aufgabe}{16}

%% TODO Titel: Ich-kann-Satz fehlt in der Bank (Platzhalter: Kettenname) – Nr. 17
%% TODO Anweisung: ein Satz über der ersten Teilaufgabe fehlt in der Bank – Nr. 17
\begin{aufgabe}{Wurzelgesetze und rationale Exponenten}
\begin{teile}
\teil $\sqrt{4 \cdot 81}$ – was steht unter der Wurzel? \\ \kreuz{Produkt oder Quotient: trennbar} \\ \kreuz{Summe oder Differenz: erst ausrechnen}
\end{teile}
%% TODO Darstellung für schwach fehlt in der Bank (reelle-zahlen-e3-k1-s1-v1; Abschnitt „Für schwache Schüler“)
\swz{$\sqrt{9 \cdot 64}$ – erst malnehmen, dann Wurzel; und die Wurzeln einzeln. Gleiches Ergebnis?}{}
%% TODO Darstellung für schwach fehlt in der Bank (reelle-zahlen-e3-k1-s1-v2; Abschnitt „Für schwache Schüler“)
\swz{$\sqrt{4 \cdot 121}$ – erst malnehmen, dann Wurzel; und die Wurzeln einzeln. Gleiches Ergebnis?}{}
%% TODO Darstellung für schwach fehlt in der Bank (reelle-zahlen-e3-k1-s1-v3; Abschnitt „Für schwache Schüler“)
\swz{$\sqrt{9 \cdot 100}$ – erst malnehmen, dann Wurzel; und die Wurzeln einzeln. Gleiches Ergebnis?}{}
%% TODO Darstellung für schwach fehlt in der Bank (reelle-zahlen-e3-k1-s1-v4; Abschnitt „Für schwache Schüler“)
\swz{$\sqrt{64 \cdot 81}$ – erst malnehmen, dann Wurzel; und die Wurzeln einzeln. Gleiches Ergebnis?}{}
%% TODO Darstellung für schwach fehlt in der Bank (reelle-zahlen-e3-k1-s1-v5; Abschnitt „Für schwache Schüler“)
\swz{$\sqrt{121 \cdot 9}$ – erst malnehmen, dann Wurzel; und die Wurzeln einzeln. Gleiches Ergebnis?}{}
\swfrage{Was bleibt gleich, was ändert sich?}
%% TODO Darstellung für schwach fehlt in der Bank (reelle-zahlen-e3-k1-s2-v1; Abschnitt „Für schwache Schüler“)
\swz{$\sqrt{\frac{64}{9}}$ – als Bruch?}{}
%% TODO Darstellung für schwach fehlt in der Bank (reelle-zahlen-e3-k1-s3-v1; Abschnitt „Für schwache Schüler“)
\swz{$\sqrt{28}$ – teilweise Wurzel ziehen?}{}
%% TODO Darstellung für schwach fehlt in der Bank (reelle-zahlen-e3-k1-s4-v1; Abschnitt „Für schwache Schüler“)
\swz{$4\sqrt{3} + 2\sqrt{3}$ – zusammengefasst?}{}
%% TODO Darstellung für schwach fehlt in der Bank (reelle-zahlen-e3-k1-s5-v1; Abschnitt „Für schwache Schüler“)
\swz{$\sqrt{81 + 144}$ – Ergebnis?}{}
%% TODO Darstellung für schwach fehlt in der Bank (reelle-zahlen-e3-k1-s6-v1; Abschnitt „Für schwache Schüler“)
\swz{$81^{\frac{1}{2}}$ – als Wurzel und als Zahl?}{}
%% TODO Darstellung für schwach fehlt in der Bank (reelle-zahlen-e3-k1-s7-v1; Abschnitt „Für schwache Schüler“)
\swz{$\sqrt[3]{64}$ – Ergebnis?}{}
%% TODO Darstellung für schwach fehlt in der Bank (reelle-zahlen-e3-k1-s8-v1; Abschnitt „Für schwache Schüler“)
\swz{$5^{\frac{1}{2}} \cdot 5^{\frac{1}{2}}$ – Ergebnis?}{}
%% TODO Darstellung für schwach fehlt in der Bank (reelle-zahlen-e3-k1-s9-v1; Abschnitt „Für schwache Schüler“)
\swz{$\frac{1}{\sqrt{3}}$ – mit rationalem Nenner?}{}
%% TODO Darstellung für schwach fehlt in der Bank (reelle-zahlen-e3-k1-s10-v1; Abschnitt „Für schwache Schüler“)
\swz[3]{Begründe mit den Potenzgesetzen: $\sqrt{a} \cdot \sqrt{b} = \sqrt{a \cdot b}$ für $a, b \ge 0$.}{}
%% TODO Darstellung für schwach fehlt in der Bank (reelle-zahlen-e3-k1-s11-v1; Abschnitt „Für schwache Schüler“)
\swz{$x^3 = 512$ – Lösung?}{}
\swz[3]{$\sqrt{18} + \sqrt{8}$ – vereinfacht? Schreibe das Ergebnis auch mit einer Potenz von $2$.}{}
\end{aufgabe}

%% TODO Titel: Ich-kann-Satz fehlt in der Bank (Platzhalter: Kettenname) – Nr. 18
%% TODO Anweisung: ein Satz über der ersten Teilaufgabe fehlt in der Bank – Nr. 18
\begin{aufgabe}{Wurzelgesetze und rationale Exponenten – Fehler finden, Begründen, Darstellungswechsel, Anwendung}
\swz[3]{Finde den Fehler und rechne richtig: \rechnung{\sqrt{81} + \sqrt{144} &= \sqrt{225} \\ &= 15}}{}
\swfrage{Zeige mit einem Zahlenbeispiel, dass $\sqrt{a + b}$ im Allgemeinen nicht $\sqrt{a} + \sqrt{b}$ ist.}
\swz{$\sqrt[5]{x^3}$ – als Potenz?}{}
\swz[3]{Ein quadratisches Grundstück hat $800\,\text{m}^2$. Gib die Seitenlänge exakt und möglichst einfach sowie gerundet an.}{}
\end{aufgabe}

\uebersichtskasten{Wurzeln malnehmen und teilen: Unter einer Wurzel darf man Produkte und Quotienten trennen: √4 · √9 = √36 = 6; √50 = √25 · √2 = 5√2 (teilweises Wurzelziehen). \\ Keine Wurzel aus Summen trennen: √(a + b) ist nicht √a + √b – erst die Summe ausrechnen: √(9 + 16) = √25 = 5, aber √9 + √16 = 3 + 4 = 7. \\ Wurzel als Potenz: Die Quadratwurzel ist die Potenz mit dem Exponenten ein halb, die n-te Wurzel hat den Exponenten 1/n – dann gelten die Potenzgesetze auch für Wurzeln: √a = $a^{(1/2)},$ ³√8 = $8^{(1/3)}$ = 2 (Taschenrechner: Klammer um den Exponenten).}
